\section{Experimental plan}
\label{sec:b-experiments}

No result in this section exists yet.  Every performance number is
\emph{planned}; the figure and table slots name the generating script and the
dataset they require, and none of them contains a fabricated curve.

\subsection{Hypotheses}

\begin{table}[t]
\centering
\small
\setlength{\tabcolsep}{4pt}
\begin{tabular}{@{}p{0.07\linewidth}p{0.52\linewidth}p{0.35\linewidth}@{}}
\toprule
ID & Hypothesis & Falsified by \\
\midrule
H1 & The learned advantage allocation beats the matched failure-risk allocation
at the primary capacity on held-out templates. & $D_O(q)$ not distinguishable
from zero at $q=0.25$. \\
H2 & The learned advantage allocation beats the plain action-conditioned value
controller at the same capacity and information. & Same contrast against
baseline (f). \\
H3 & Conservative selection retains a positive mean gain while limiting
no-injected-fault degradation relative to always-verify. & Gain collapses, or
no-fault degradation exceeds the predeclared tolerance at every qualifying
operating point. \\
H4 & The interaction of feasibility, cost and harm changes a meaningful
decision: capacity spent on unavailable checkpoints, or on checkpoints whose
verification is harmful, changes the achieved gain. & The no-feasibility or
no-harm ablations leave the allocation unchanged in the qualified pool. \\
\bottomrule
\end{tabular}
\caption{Planned hypotheses and their falsifiers.  H1--H3 are the empirical
core; H4 is a focused structural check, not a sweep.}
\label{tab:b-hypotheses}
\end{table}

\subsection{Studies}

\begin{table}[t]
\centering
\small
\setlength{\tabcolsep}{4pt}
\begin{tabular}{@{}p{0.06\linewidth}p{0.40\linewidth}p{0.44\linewidth}@{}}
\toprule
ID & Question & Data and method \\
\midrule
B0 & Can the declared live checkpoint be branched consistently? & B-Q pool and
controlled fixture checks; runtime boundary and qualification report. \\
B1 & Is the task/intervention pair scientifically and computationally usable? &
Disjoint pilot, fixed mechanics, bounded candidate count; yield, feasibility,
variance, effect uncertainty, runtime and scale estimate. \\
B2 & Do learned advantage allocations improve over strong matched alternatives? &
Held-out branch data and matched budgets; policy value differences, support,
intervals and allocation curves. \\
B3 & What is gained and lost by conservative selection? & Frozen policy family,
independent calibration/test; utility, invalidity, no-injected-fault
performance, cost and intervention rate. \\
B4 & Which structural assumptions matter? & Small predeclared ablations: mask
removed, cost ignored, harm ignored. \\
\bottomrule
\end{tabular}
\caption{First-study sequence.  B5 (recovery) and B6 (second deployment
replication) are optional and not on the critical path.}
\label{tab:b-studies}
\end{table}

\subsection{Pilot envelope and decision points}

The planning envelope is $20$--$40$ independent development templates, one
checkpoint per root, $R=2$--$4$ repetitions, if compute permits; it is a
development suggestion, not a powered sample-size claim.  The pilot verifies
action semantics and endpoint validity, estimates runtime per
prefix/branch/adjudication, measures feasibility and checkpoint reachability,
estimates paired outcome variation and template clustering, checks whether
intervention effects are learnable with available features, and estimates
final-test precision for the actual paired policy contrast.  The smallest
scientifically useful gain $\epsilon$, a harm tolerance and a resource ceiling
are recorded before final-test outcomes and are not weakened afterwards.  If
the planned precision is unaffordable, claims are narrowed or the study is
changed prospectively; that is not disguised as the original confirmatory run.
Counts are reported separately as templates, roots, checkpoints, repetitions
and executed arms.

\subsection{Required figures and tables (slots)}

\begin{table}[t]
\centering
\small
\setlength{\tabcolsep}{4pt}
\begin{tabular}{@{}p{0.24\linewidth}p{0.34\linewidth}p{0.34\linewidth}@{}}
\toprule
Artifact & Scientific role & Required data / generator \\
\midrule
Risk/deficit versus advantage & Actionability patterns with uncertainty &
Out-of-sample predictions plus grouped held-out outcome summaries;
\texttt{experiments/paper\_b/figures/risk\_vs\_advantage.py} \\
Gain versus budget & Primary matched allocation comparison &
Frozen policies, paired evaluation and uncertainty;
\texttt{experiments/paper\_b/figures/gain\_vs\_budget.py} \\
Utility--harm--intervention tradeoff & What conservative selection buys and
costs & Calibration-defined operating points and held-out metrics;
\texttt{experiments/paper\_b/figures/tradeoff.py} \\
Availability/cost interaction & Why the combined decision problem matters &
Qualified masks and predeclared cost conditions;
\texttt{experiments/paper\_b/figures/availability\_cost.py} \\
Dataset/support table & Effective sample sizes and exclusions & All roots
including unreachable, unavailable, failed and invalid;
\texttt{experiments/paper\_b/figures/support\_table.py} \\
Runtime/resource table & Feasibility and overhead & Prefix, branch, verifier,
snapshot and adjudication costs;
\texttt{experiments/paper\_b/figures/runtime\_table.py} \\
\bottomrule
\end{tabular}
\caption{Figure and table slots for the first study.  Each generator is a slot
that must read a qualified dataset; none is populated with invented numbers.}
\label{tab:b-figure-slots}
\end{table}

\subsection{What a rejection does and does not establish}
\label{sec:b-rejection}

A positive held-out $D_O(q)$ demonstrates superiority of the two implemented
policies under the studied law.  It does \emph{not} by itself reject the
proposition that true outcome risk and true advantage have identical
orderings: estimation error in the learned risk model can make it suboptimal
even under that population null.  The primary null is therefore stated as a
comparison of the frozen learned policies; risk calibration and discrimination
are reported with matched model capacity; and the exact ordinal-sufficiency
theorem is reserved for population quantities and known-truth constructions.
Non-rejection is inconclusive: it does not establish equivalent orderings or the
absence of actionable heterogeneity.
